% !PRINTED 2
\weberpage{2}
%<<CONT
\gesperrt{number}. The number $1$ itself is customarily not designated
a prime number, for reasons of expediency. If two numbers are divisible
by a third, then the sum and the difference of the first two are also
divisible by the third; and if one number is divisible by a second and
this second by a third, then the first is also divisible by the third.

Two numbers always have the common divisor $1$. If they have no other
common divisor, as for example the pairs of numbers $5$ and $7$, or
$21$ and $38$, then the two numbers are called \gesperrt{relative prime
numbers}, or also \gesperrt{numbers prime to one another}.

Among the common divisors of any two given numbers one will be the
greatest, and it is a very important problem to find this greatest
common divisor. This is accomplished by a procedure known under the
name \gesperrt{algorithm of the greatest common divisor}, which already
occurs in Euclid.\footnote{Elements, Book VII, II, vol. II of Heiberg's
edition.}% Source prints Roman II (proposition 2), not Arabic 11.

Let $a,a_1$ be the two numbers; take the larger of them, say $a$, as
dividend and the smaller $a_1$ as divisor, and determine the quotient
$q_1$ and, if the division does not come out exactly, the remainder
$a_2$, so that
\[
 a=q_1a_1+a_2,
\]
and $a_2<a_1$. Every common divisor of $a$ and $a_1$ is then also a
common divisor of $a_1$ and $a_2$, and conversely. If one proceeds with
$a_1,a_2$ just as with $a,a_1$, and puts, when the division does not
come out exactly,
\[
 a_1=q_2a_2+a_3,
\]
then again $a_3<a_2$, and every common divisor of $a_1$ and $a_2$ is
also a common divisor of $a_2$ and $a_3$, and conversely. If one
continues dividing in this way, then, since the numbers
$a,a_1,a_2,a_3,\ldots$ continually decrease, the division must
necessarily come out exactly after a finite number of steps; that is,
at last a pair of equations
\[
 a_{\nu-2}=q_{\nu-1}a_{\nu-1}+a_\nu,
 \qquad
 a_{\nu-1}=q_\nu a_\nu
\]
must occur. Then $a_\nu$ is a common divisor of all preceding $a$'s,
hence also of $a$ and $a_1$, and every common divisor of $a$ and $a_1$
is a divisor of $a_\nu$. Thus $a_\nu$ is the greatest common divisor of
$a$ and $a_1$, and at the same time we have arrived at the theorem that
every common divisor of two numbers must divide their greatest common
divisor.
